\section{Synthesis: From Evidence to Reuse Decisions}\label{sec:discussion}
\begin{table*}[t]\centering\footnotesize
\caption{Claim-conditioned reuse decisions derived from the record and holdout audits. Each changed boundary identifies the evidence that must be bound before carrying a result forward.}\label{tab:r19-decisions}
\begin{tabularx}{\textwidth}{@{}p{.17\textwidth}p{.29\textwidth}Y@{}}\toprule
Reuse question & Evidence decision & Retain or measure next\\\midrule
Emitter / representation & An upstream or buffered capture can differ from the final receiving trace and delivery schedule. & Bind inference to the final receiver; retain the emitter's overflow, byte and delivery-delay scope (R2/R3/R9; \S\ref{sec:discussion}).\\
Read cover $\rightarrow$ writes & Discarding a reply does not establish unchanged commits, authorization or workflow outcomes. & Bind the protected state and utility to the executed intervention; include cover calls, quotas and failures (R3/R8/R9).\\
Customer / model calibration & A reported split can separate inputs without changing the generator or calibration population. & State held-out units and calibration provenance; separate local fitting variability from source/domain changes; evaluate the requested transfer axis (R5/R6; \S\ref{sec:evaluation}).\\
Attacker observation & An outcome for one projection or fitted attacker does not supply a result for a richer receiving view. & Evaluate comparable defended tests under the intended observation and adaptation contract; retain receiver-specific inference metrics (R2/R5/R7).\\
Stored trace $\rightarrow$ online path & A trace transformation or simulated span does not measure executed queues, completion or online latency. & Bind the receiver, harness and resource denominator to the online condition; establish the relevant topology/network holdout (R3/R6/R9).\\
Bound cache-reuse package & Early Bird prefix recovery and CachePrune direct recovery permit a scoped candidate comparison; CachePrune binds a shared utility/cost package. & Retain that package for its measured configuration; measure prefix-policy serving cost separately, and re-evaluate any changed workload/hardware.\\
Component $\rightarrow$ workflow utility & Component quality or logical state conformance does not determine a planner's response to scheduling and deadlines. & Report utility at the requested workflow boundary and the same intervention/configuration as protection and cost (R3/R8/R9).\\\bottomrule
\end{tabularx}\end{table*}

The census (\S\ref{sec:defenses}) covers all 39 records and the holdout matrix (\S\ref{sec:evaluation}) all 82 families. We now select three examples by the \emph{comparison decision} they expose, not by effect size: a changed receiving path, cover with different state effects, and an adaptation/cost mismatch within a named mechanism. Their value is the decision they change: retain an aligned result, separate a wrongly joined package, or require a result at a changed execution boundary.

\subsection{Case 1: bind the final target-bearing emitter}
NetEcho's upstream packet-size shaping operates on a different projection from the tool-call sequence. Its plaintext-aware frontend then emits a new size trace (grouped token counts or character lengths), reopening a size channel after upstream protection \cite{netecho2025}. SplitDecoy is a separate return-path record: the cloud sees nonzero gradient row support on every tested frame across nine seeds \cite[Sections 5.1--5.3]{politis2026splitdecoy}. Both change the receiver boundary and therefore require a separately bound inference result.

The data processing inequality requires $S\to Z\to Y$: target-oblivious postprocessing cannot increase $I(S;Z)$ \cite[Theorem~2.8.1]{cover2006elements}. Plaintext-aware re-emission or backward computation breaks that premise. The literature comparison must therefore bind both the final receiving projection and the emitter whose costs are measured.

\subsection{Case 2: resolve the WRA versus AHB cover asymmetry}
WRA's one-decoy $E_{\rm func}$ drops from 0.770 to 0.610, whereas AHB's Trip Planning success falls from 0.75 to 0.15 under \texttt{Online+Fake} humanization \cite{jeong2026wra,zhu2026turing}. Grouping both as ``cover traffic'' conceals different objectives and executions.

WRA discards decoy reads for Type~I confidentiality; AHB executes GUI events for Type~III evasion, potentially navigating or submitting forms. Their utility outcomes are not contradictory, and AHB's missing condition-specific denominators preclude a matched causal effect. A shared workflow setting does not establish shared state effects. Terminal-state reward and dispatch-time authorization also answer different questions: an unconfirmed write can still reach the intended final state. Reuse requires an explicit state/authorization outcome at the intervention being carried forward.

\subsection{Case 3: bind adaptive retraining and cost objects}
A matched adaptation contrast compares clean-trained $D_0$ and defense-trained $D_2$ on identical defended tests. UI Traces and Dangol raise the record-binding question this leaves open: whether a refit, a measured delay and a utility outcome refer to the same executed condition \cite{uitraces2026,dangol2026workflow}. Cost objects show the same pattern. Early Bird's $6$ms/$4\%$ GPTCache cost does not price $K$-prefix isolation; KVGov's $1.3\,\mu$s HMAC microcost omits lost reuse; SpliceDefense's QCP cost does not price CTBF+QCP \cite{earlybird2025,addagada2026kvgov,spliceleak2026}. CachePrune binds protection, QA/summarization utility and throughput, while OTRO separates conditional leakage from TTFT/padding costs \cite{wu2026cacheprune,lee2026otro}. The reuse decision follows the evaluated configuration: retain each source's bound outcome, and request the missing adaptation, execution or cost measurement before extending it.


\subsection{A positive reuse path: from comparison to a candidate configuration}
The Early Bird $K$-prefix record and CachePrune's direct-recovery record form one of the 17 descriptive matches: both address sensitive-token recovery through cache-reuse probes \cite[V-A]{earlybird2025}, \cite[5.1--5.6]{wu2026cacheprune}. This permits an analyst to retain both mechanisms as candidates for that threat, while keeping the oracle and executed-serving measurements separate. The joint audit then changes the next step. CachePrune supplies a default configuration package linking recovery, generation quality and serving cost; the analyst can use that package for a source-scoped feasibility assessment. Early Bird's semantic-cache overhead cannot price its $K$-prefix candidate, so that candidate still needs a serving-cost measurement. Neither a pooled recovery ranking nor a new-hardware guarantee follows from this comparison.

\paragraph{What the three findings add}
Record-level comparison turns a mechanism-level shortlist into correctly scoped alternatives. Configuration-level joint binding distinguishes a reusable package from borrowed utility or cost. Deployment decisions then identify the changed object and the measurement it needs. Compared with established WF practice on adaptive attacks, shifts and online costs \cite{juarez2014critical,mathews2023critical}, the LLM/agent synthesis makes three additional joins explicit: prefix versus semantic-cache primitives, discarded read outputs versus executed state changes, and component generation quality versus planner outcomes. The contribution is the source-backed decision change, not a new information-processing theorem.

\subsection{A deployment query determines the evidence to retain}
Start with the intended receiver, protected outcome and intervention. Record-level comparison identifies whether the cited result measures that object. Configuration-level joint binding identifies which inference, utility and cost outcomes share its configuration. Deployment decisions use the holdout matrix to identify evidence for the proposed change and specify the remaining measurement. Table~\ref{tab:r19-decisions} makes these outputs operational: a new emitter needs a receiver-specific capture, writes need state/authorization checks, source-calibration comparisons need a local fitting-variability control, and an online scheduler needs its own planner utility and latency measurement. The calibration control follows from the comparison design: changing the source can also change the finite training sample used to fit the defense.

Appendix~\ref{app:case-study} applies this procedure in one compact case study. The literature audit and the decision table establish the SoK contribution; the case study provides a worked application with separately scoped measurements.
